\bookchapter{Kapitel 4: Apps installieren und verwalten – der „Fähigkeitsrucksack“ des Smartphones}

\chapterintro{Apps machen das Smartphone zu einem Allrounder, aber sie müssen auch verwaltet werden. Genau wie du einem neuen Mitarbeiter einen Arbeitsplatz zuweist und Berechtigungen vergibst, erfordert die Installation von Apps einen standardisierten Prozess.}

\booksection{Eins. Offizielle Stores: Der sicherste „App-Supermarkt“}

\begin{itemize}
  \item \textbf{Android}: Google Play (international) oder markeneigene Stores (z. B. Huawei AppGallery, Xiaomi App Store).
  \item \textbf{iOS}: Nur der App Store.
  \item \textbf{HarmonyOS}: Huawei AppGallery.
\end{itemize}

Zentrale Sicherheitsbedeutung: geprüft, Updates garantiert, sichere Zahlungen. Handlungsleitfaden:\textbf{Downloads vorzugsweise nur aus offiziellen Stores!}

\booksection{Zwei. Reiner Modus / Schutzmodus: Den „Virenschutz“ des Smartphones aktivieren}

Alle großen chinesischen Marken (Huawei, Xiaomi, OPPO, vivo, Honor) verfügen über einen integrierten\textbf{reinen Modus}bzw.\textbf{bzw.}Schutzmodus\textbf{– einen systemweiten Sicherheitsschalter.}Android

\booksubsection{Was ist das eigentlich?}

iOS\textbf{Blockieren der Installation von Drittanbieter-Apps}HarmonyOS\textbf{erlaubt das System nur Installationen aus dem offiziellen App-Store}. Installationspakete, die über Browser, QR-Codes, Chat-Software usw. heruntergeladen werden, werden direkt abgefangen – Malware wird an der Quelle unterbunden.

\booksubsection{Zielgruppe}

\begin{itemize}
  \item \textbf{Ältere Nutzer}: Die „offiziell" von „gefälscht" nicht unterscheiden können – der reine Modus ist der zuverlässigste Schutz.
  \item \textbf{Nutzer, die App-Berechtigungen nicht kritisch sehen}: Keine kleinen App-Stores, Testversionen oder ausländische Spiele nötig haben.
  \item \textbf{Alle Nutzer, die die Gerätesicherheit erhöhen möchten}: Selbst als Power-User ist Aktivieren besser als Deaktivieren – bei tatsächlichem Bedarf einfach kurz ausschalten.
\end{itemize}

\booksubsection{Aktivierungspfad (Beispiel Huawei/Xiaomi)}

\begin{itemize}
  \item \textbf{.}: [Einstellungen] > [System und Updates] >\textbf{[Reiner Modus]}, „Einschalten" wählen.
  \item \textbf{Xiaomi}: [Einstellungen] > [Datenschutz] >\textbf{[Datenschutz schützen]}oder [Einstellungen] > [App-Einstellungen] > [App-Installation], „Installation unbekannter Apps" deaktivieren.
  \item \textbf{OPPO/vivo}: [Einstellungen] > [Sicherheit] >\textbf{[Installation von Apps aus externen Quellen]}, Berechtigungen einzeln entziehen.
\end{itemize}

\booksubsection{Sicherheitsprüfmechanismus der App-Stores}

Ein App-Store ist keine „Dekoration“ – hinter ihm steht eine vollständige Sicherheits-Pipeline:

\begin{itemize}
  \item \textbf{Virenscan}: Vor der Aufnahme jeder App wird ein gemeinsamer Scan durch mehrere Antiviren-Engines durchgeführt, um bekannte Trojaner und Virensignaturen zu erkennen.
  \item \textbf{Berechtigungsprüfung}: Der Store prüft, ob die von der App angeforderten Berechtigungen „plausibel" sind – eine Taschenlampen-App möchte auf das Adressbuch zugreifen? Direkt abgelehnt oder entfernt.
  \item \textbf{Entwicklerzertifizierung}: Für den Store-Einstieg braucht es Firmenqualifikation, Personalausweis- oder Kreditkartenprüfung. Bei Problemen gibt es einen Ansprechpartner – keine „Schwarz-App“.
  \item \textbf{Kontinuierliche Überwachung}: Auch gelistete Apps werden regelmäßig neu gescannt; neu entdecktes Schadverhalten führt zur erzwungenen Löschung.
\end{itemize}

\booksection{Drittanbieter-Quellen und „elektronische Parasiten“: höchst riskant!}

Apps aus Webseiten oder unbekannten Links können Viren einschleppen, Privatsphäre ausspionieren und Zahlungsdaten abgreifen. Ein neues Konzept:\textbf{„elektronische Parasiten“}.

\textbf{Was sind „elektronische Parasiten“?}\par

Gemeint sind\textbf{„geknackte“ Software mit eingebautem Trojaner}– du glaubst, Geld zu sparen, und bist längst reingefallen. Typische Formen:

\begin{itemize}
  \item „Crack-Versionen“, „Gratis-Versionen“, „werbefreie Versionen“ kostenpflichtiger Software;
  \item graue Tools wie „Likes klicken“, „Rote-Umschläge-Abgreifen“, „Spiel-Cheats“;
  \item „internen Versionen“, die in Gruppenchats per Kleinanzeige empfohlen werden.
\end{itemize}

Diese Programme funktionieren an der Oberfläche – im Hintergrund können sie\textbf{Kontakte stehlen, SMS lesen, das Mikrofon abhören und im Hintergrund Coins schürfen}. Der reine Modus ist der\textbf{erste Verteidigungswall}– er verhindert die Installation einfach komplett.

\booksubsection{Empfehlung für den zweiten Verteidigungswall:}

\begin{itemize}
  \item \textbf{Regelmäßiger kompletter Virenscan}: Nutze das integrierte Sicherheitscenter (z. B. Huawei-Telefonmanager, Xiaomi-Sicherheitscenter) einmal pro Monat.
  \item \textbf{Sensible Berechtigungen vorsichtig vergeben}: Kontakte, SMS, Standort, Kamera, Mikrofon – diese fünf „Hochrisiko-Berechtigungen“ streng prüfen.
  \item \textbf{Verbrauch und Datenvolumen im Blick}: Plötzlich heißes Handy, schnell leerer Akku, nächtliches Datenvolumen – zu 80 \% steckt ein „Parasit“ im Hintergrund dahinter.
\end{itemize}

Falls es wirklich nicht ohne eine Drittanbieter-App geht:

\begin{enumerate}
  \item Stelle absolut sicher, dass die Quelle vertrauenswürdig ist.
  \item Prüfe vor der Installation die angeforderten Berechtigungen genau.
  \item Scanne danach mit einer Sicherheits-App.
\end{enumerate}

Letzter Rat: In den Systemeinstellungen die Berechtigung „Apps aus unbekannten Quellen installieren“ ausschalten.

\booksection{Installieren, aktualisieren, deinstallieren}

\begin{itemize}
  \item \textbf{Installieren}: Im Store „Abrufen“ oder „Installieren“ tippen; beim ersten Start fragt die App nach ihren Berechtigungen.
  \item \textbf{Aktualisieren}:\textbf{Unbedingt updaten!}Vor allem werden Sicherheitslücken gestopft. Aktiviere die automatische Aktualisierung.
  \item \textbf{Deinstallieren}: Icon lange drücken und in den Deinstallationsbereich ziehen. Für eine gründlichere Räumung: [Einstellungen] > [App-Verwaltung] – dort finden und deinstallieren.
\end{itemize}

\booksection{Berechtigungsmanagement: du bist der „Chef“ der Apps}

Apps beantragen Rechte (Kontakte, Standort) mal zu Recht (Karten-App braucht Standort), mal zu Unrecht (die Taschenlampe will die Kontakte).

\begin{itemize}
  \item \textbf{Berechtigungsoptionen}: „Immer erlauben“, „Nur bei Verwendung erlauben“, „Nicht erlauben“, „Jedes Mal fragen“.
  \item \textbf{Leitprinzipien}:
  \begin{enumerate}
    \item Minimierung: nur geben, wenn nötig.
    \item Nachkontrolle: regelmäßig in den Berechtigungseinstellungen sehen, wer „über seine Rechte“ hinausgeht.
    \item Situationabhängig: Der Taxi-App den Standort nur während der Nutzung geben.
  \end{enumerate}
\end{itemize}

\begin{tcbwarn}{Berechtigungen mit Fokus überwachen}
\begin{itemize}
  \item Kontakte/SMS: könnten für Spam-SMS missbraucht werden.
  \item Echtzeit-Standort: könnte deine Bewegungen verfolgen.
  \item Kamera/Mikrofon: könnten heimlich filmen und aufnehmen.
  \item Dateizugriff: könnte deine privaten Dateien lesen.
\end{itemize}

\end{tcbwarn}

\booksection{Aufbau eines Software-Installationspakets}

: Der „Ausweis“ und die „Bedienungsanleitung” der App. Verzeichnet den Paketnamen, die Versionsnummer, benötigte Berechtigungen, unterstützte Systemversionen, Einstiegs-Activity und weitere Kerninformationen. Bei der Installation liest das System diese Datei zuerst.

\booksubsection{Ein APK (Android Package Kit) ist im Wesentlichen ein ZIP-Archiv. Nach dem Entpacken siehst du folgende Kernverzeichnisse und Dateien:}

Eine APK (Android Package Kit) ist im Kern ein ZIP-Archiv. Nach dem Entpacken sieht man folgende Kernverzeichnisse und Dateien:

\begin{itemize}
  \item \textbf{AndroidManifest.xml}: Der „Ausweis“ und die „Betriebsanleitung“ der App. Enthält Paketname, Version, benötigte Berechtigungen, unterstützte Systemversionen, Einstiegs-Activity und weitere Kerndaten. Bei der Installation liest das System diese Datei zuerst.
  \item \textbf{:}: Der Kerncode der App. Java-/Kotlin-Quellcode wird zu Dalvik-Bytecode kompiliert und in dieser Datei gespeichert. Eine App kann mehrere dex-Dateien enthalten (z. B. classes2.dex, classes3.dex).
  \item \textbf{resources.arsc}: Die kompilierte Ressourcen-Indextabelle. Verzeichnet die IDs aller Ressourcen und ihre Zuordnungen.
  \item \textbf{res/-Verzeichnis}: Ressourcenordner. Enthält Bilder, Layout-Dateien, Zeichenketten, Farbdefinitionen usw. Bilder unterschiedlicher Auflösung liegen in verschiedenen drawable-Ordnern (z. B. drawable-hdpi, drawable-xhdpi).
  \item \textbf{assets/-Verzeichnis}: Ordner für Rohressourcen. Enthält nicht kompilierte Ressourcen wie Webseiten-Dateien, Schriftarten, Audiodateien. Zur Laufzeit kann die App diese über den AssetManager lesen.
  \item \textbf{lib/-Verzeichnis}: Ordner für native Bibliotheken. Enthält kompilierte .so-Dateien (C/C++-Kompilat), nach CPU-Architektur sortiert (z. B. arm64-v8a, armeabi-v7a). Spiele und rechenintensive Apps verfügen in der Regel über dieses Verzeichnis.
  \item \textbf{META-INF/-Verzeichnis}: Ordner mit Signaturinformationen. Enthält die drei Dateien MANIFEST.MF, CERT.SF und CERT.RSA zur Verifizierung der Integrität und Signatur des APK.
\end{itemize}

\booksubsection{Ein IPA (iOS App Store Package) ist ebenfalls ein komprimiertes Archiv, dessen Struktur jedoch leicht abweicht:}

: Die Konfigurationsdatei der App, ähnlich der AndroidManifest.xml bei Android. Verzeichnet die App-ID, Version, Berechtigungsanforderungen, unterstützte Geräte usw.

\begin{itemize}
  \item \textbf{Payload/-Verzeichnis}: Hier liegt der Kerninhalt. Darin befindet sich ein .app-Ordner, der den gesamten App-Inhalt enthält.
  \item \textbf{Info.plist}: Die Konfigurationsdatei der App, vergleichbar mit der AndroidManifest.xml unter Android. Enthält App-ID, Version, Berechtigungsbedarf, unterstützte Geräte und mehr.
  \item \textbf{Ausführbare Datei}: Das Hauptprogramm, das genauso heißt wie der .app-Ordner – eine kompilierte Mach-O-Binärdatei.
  \item \textbf{Frameworks/-Verzeichnis}: Dynamische Bibliotheken und Frameworks, von denen die App abhängt.
  \item \textbf{Ressourcendateien}: Bilder, Audios, Videos und andere Ressourcen, normalerweise im .app-Ordner gespeichert.
  \item \textbf{\_CodeSignature/-Verzeichnis}: Code-Signaturinformationen, mit denen sich Integrität und Herkunft der App prüfen lassen.
\end{itemize}

\booksubsection{Der Signaturprüfmechanismus}

: Dynamische Bibliotheken und Frameworks, von denen die App abhängt.

\begin{itemize}
  \item \textbf{Android-Signatur}: Der Entwickler signiert das APK mit einem privaten Schlüssel; bei der Installation verifiziert das System mit dem öffentlichen Schlüssel. Wurde das APK manipuliert (z. B. mit einem Virus infiziert), wird die Signatur ungültig und das System verweigert die Installation. Google Play verlangt zudem App Signing, bei dem Google den Signaturschlüssel verwahrt und die App-Sicherheit zusätzlich schützt.
  \item \textbf{iOS-Signatur}: Apples Prüfung ist noch strenger. Jedes IPA muss mit einem von Apple ausgestellten Zertifikat signiert sein, sonst lässt es sich nicht installieren. Deshalb können iOS-Apps nur aus dem App Store oder über Unternehmenszertifikate installiert werden. Apps mit Unternehmenszertifikat haben eine begrenzte Gültigkeit und Apple kann sie jederzeit widerrufen.
  \item \textbf{Wozu die Signatur dient}:
  \begin{enumerate}
    \item Herkunft prüfen: bestätigt, dass die App wirklich vom angegebenen Entwickler stammt.
    \item Integrität prüfen: stellt sicher, dass die App unterwegs nicht manipuliert wurde.
    \item Fälschung verhindern: Eine Schadsoftware kann sich nicht als andere App ausgeben.
  \end{enumerate}
\end{itemize}

\booksubsection{Paketgröße und -optimierung}

Quelle verifizieren: Bestätigen, dass die App tatsächlich vom angeblichen Entwickler stammt.

\begin{itemize}
  \item \textbf{Download-Größe}: Die Größe des komprimierten APK/IPA.
  \item \textbf{Installationsgröße}: Die Größe nach dem Entpacken – meist deutlich größer als die Download-Größe.
  \item \textbf{Platzbedarf im Betrieb}Integrität verifizieren: Sicherstellen, dass die App während der Übertragung nicht manipuliert wurde.
\end{itemize}

: Zur Laufzeit erzeugt die App zusätzlich Caches und Datendateien und belegt noch mehr Speicher.

\begin{itemize}
  \item \textbf{App Bundle (Android)}: Googles Play Dynamic Delivery. Beim Download erhält der Nutzer nur die für sein Gerät passenden Ressourcen (z. B. nur die Bilder zur passenden Bildschirmauflösung), was die Download-Größe deutlich reduziert.
  \item \textbf{Praxis: Berechtigungen einer App prüfen}: Ähnliche Technik – Apple schneidet nicht benötigte Ressourcen automatisch geräteabhängig weg.
  \item \textbf{Download bei Bedarf}: Große Apps wie Spiele laden zuerst den Kern und ziehen Ressourcenpakete später nach.
\end{itemize}

\begin{tcbinfo}{Wähle eine häufig verwendete App (z. B. Social Media, Shopping, Werkzeug), öffne die App-Berechtigungsverwaltung in [Einstellungen], prüfe jeden Eintrag und deaktiviere unnötige sensible Berechtigungen oder ändere sie zu „Nur bei Verwendung erlauben".}
Kapitel 5: Das Smartphone besser verstehen lassen – Grundeinstellungen

\end{tcbinfo}

\begin{tcbpractice}{Ein paar Minuten Einstellungen verbessern das Erlebnis enorm. Diese Einstellungen sind wie das „Stimmen" des Smartphones, damit es besser zu deinen Gewohnheiten passt.}
Eins. Ton und Vibration

\end{tcbpractice}
